%%%%%%%% ICML 2026 EXAMPLE LATEX SUBMISSION FILE %%%%%%%%%%%%%%%%%

\documentclass{article}

% Recommended, but optional, packages for figures and better typesetting:
\usepackage{microtype}
\usepackage{graphicx}
\usepackage{subcaption}
\usepackage{booktabs} % for professional tables
\usepackage{hyperref}

% Attempt to make hyperref and algorithmic work together better:
\newcommand{\theHalgorithm}{\arabic{algorithm}}

% Use the following line for the initial blind version submitted for review:
\usepackage[accepted]{icml2026}

\usepackage{amsmath}
\usepackage{amssymb}
\usepackage{mathtools}
\usepackage{amsthm}
\usepackage[capitalize,noabbrev]{cleveref}

%%%%%%%%%%%%%%%%%%%%%%%%%%%%%%%%
% THEOREMS
%%%%%%%%%%%%%%%%%%%%%%%%%%%%%%%%
\theoremstyle{plain}
\newtheorem{theorem}{Theorem}[section]
\newtheorem{proposition}[theorem]{Proposition}
\newtheorem{lemma}[theorem]{Lemma}
\newtheorem{corollary}[theorem]{Corollary}
\theoremstyle{definition}
\newtheorem{definition}[theorem]{Definition}
\newtheorem{assumption}[theorem]{Assumption}
\theoremstyle{remark}
\newtheorem{remark}[theorem]{Remark}

\icmltitlerunning{The Fine-Tuning Confounder in Model Merging}

\setlength{\abovedisplayskip}{2pt}
\setlength{\belowdisplayskip}{2pt}
\setlength{\abovedisplayshortskip}{0pt}
\setlength{\belowdisplayshortskip}{0pt}

\setlength{\textfloatsep}{5pt plus 1pt minus 3pt}
\setlength{\floatsep}{4pt plus 1pt minus 1pt}
\setlength{\dbltextfloatsep}{5pt plus 1pt minus 3pt}
\setlength{\dblfloatsep}{4pt plus 1pt minus 1pt}
\setlength{\intextsep}{4pt plus 1pt minus 1pt}

\begin{document}

\twocolumn[
\icmltitle{The Fine-Tuning Confounder: A Methodological Deconstruction of Representation Collapse in Multi-Task Model Merging}

\icmlsetsymbol{equal}{*}

\begin{icmlauthorlist}
\icmlauthor{The Methodologist Research Agent}{equal,gemini}
\end{icmlauthorlist}

\icmlaffiliation{gemini}{Autonomous ML Research Division, Gemini CLI Laboratory. Correspondence to: \href{mailto:methodologist@gemini-cli.org}{methodologist@gemini-cli.org}}

\icmlkeywords{Model Merging, Representation Collapse, Fine-Tuning Confounder, Activation Calibration}

\vskip 0.3in
]

\printAffiliationsAndNotice{\icmlEqualContribution}

\begin{abstract}
Multi-task model merging is a cost-effective, training-free paradigm to consolidate task-specific expert neural networks into a single, unified model. However, parameter-level interference often triggers severe ``variance collapse'' and ``representation collapse'' in deep layers of the merged backbone. While recent post-merge activation calibration techniques successfully restore performance, they typically rely on test-time operations and online inference hooks. In this work, we adopt the perspective of \textit{The Methodologist} to challenge the hidden assumptions of this literature. We identify a critical, overlooked confounding variable: the unregularized nature of expert fine-tuning. By allowing fine-tuned weights to drift excessively from their shared pre-trained initialization, standard practices introduce massive parameter-space misalignment that is the root cause of subsequent representation collapse. We systematically evaluate four fine-tuning regularization scenarios---varying weight decay and utilizing L2-SP penalties targeting the pre-trained progenitor---paired with three post-merge calibration methods. We show that while training-time regularization significantly aligns expert update directions (raising cross-task update cosine similarities from $0.04$ to $0.20$) and improves uncalibrated merged model accuracy, post-hoc activation-space calibration acts as a powerful equalizer. Specifically, our corrected Hybrid calibration (SP-TAAC + SLR-WBC) consistently restores merged model performance (achieving $\sim 43.6\%$ average accuracy, representing a $+17.4\%$ absolute increase) regardless of the training-time regularization choice. Our rigorous methodological analysis exposes the limits of training-time alignment and highlights the robustness of post-hoc representation alignment, establishing a new foundation for fair and comprehensive benchmarking in model merging.
\end{abstract}

\section{Introduction}
The pretrain-then-finetune paradigm has revolutionized deep learning, enabling highly capable base models to be specialized across diverse downstream tasks \cite{Vaswani2017, Caruana1997, Donahue2014}. However, serving dozens of these task-specific expert networks simultaneously introduces significant storage, VRAM, and context-switching overheads in production \cite{Ruder2017, Collobert2008}. To resolve this deployment bottleneck, multi-task model merging has recently emerged as an elegant, training-free alternative \cite{Wortsman2022, Choshen2022, Tang2023}. By directly interpolating the weight matrices of multiple expert models (which share a common pre-trained initialization progenitor), model merging consolidates multiple expert capabilities into a single cohesive network with zero additional parameter storage or backpropagation \cite{Ilharco2023, Jin2023b, Matena2022}.

Despite its theoretical and practical appeal, directly averaging the parameter weights of multiple experts (i.e., simple Weight Averaging) typically results in severe performance degradation. This degradation stems from intense parameter-level interference, which manifests as an exponential decay of activation variance in deep layers of the merged backbone---a phenomenon known as ``variance collapse'' or ``representation collapse'' \cite{Jordan2023, Sengupta2023, Mu2024}. As activations propagate through successive non-linear layers, variance collapse renders deeper layers incapable of extracting discriminative features, leading to catastrophic multi-task performance degradation.

To heal this variance collapse, recent paradigms have introduced intermediate representation and activation calibration, such as Sparsity-Preserving Task-Agnostic Activation Calibration (SP-TAAC) \cite{Anonymous2023}, REPAIR \cite{Jordan2023}, and SVD-based Low-Rank Weight and BatchNorm Calibration (SLR-WBC) \cite{Anonymous2024c}. While these methods successfully restore performance, they typically rely on post-hoc online activation hooks or offline weight updates computed under the assumption that the baseline merged model is fundamentally and inevitably broken.

In this work, we adopt the perspective of \textbf{The Methodologist} to critically re-evaluate this assumption. We argue that representation collapse is not an inevitable physical law of weight averaging, but rather a confounding consequence of unconstrained, unregularized expert fine-tuning. Standard practices fine-tune each expert independently with standard SGD and minimal weight decay. By allowing the models' weights to drift excessively from their shared pre-trained progenitor, these practices introduce severe, unnecessary parameter-space misalignment.

To investigate this hypothesis, we set up a highly rigorous, controlled evaluation pipeline. We train ResNet-18 experts on three distinct datasets (MNIST, Fashion-MNIST, and CIFAR-10) under four different training-time regularization scenarios, varying weight decay and L2-SP penalties. We then evaluate these experts under three post-merge calibration configurations (No Calibration, SP-TAAC, and Hybrid SP-TAAC + SLR-WBC). Our empirical findings reveal a highly nuanced and insightful picture of the model merging landscape:

\begin{itemize}\itemsep=0pt \parsep=0pt
    \item \textbf{Regularization Aligns Trajectories}: High training-time regularization (weight decay = $1\times 10^{-2}$) significantly constrains parameter drift, causing independent expert update directions to align. This raises cross-task update cosine similarities from a nearly orthogonal $0.04$ to over $0.20$.
    \item \textbf{Suppression of Collapse}: Constraining parameter drift directly suppresses representation collapse, leading to steady improvements in the uncalibrated merged model's average accuracy (from $26.20\%$ to $27.88\%$).
    \item \textbf{The Post-Hoc Equalizer}: Despite the clear benefit of regularization on the uncalibrated baseline, post-merge activation-space calibration acts as a powerful equalizer. When Hybrid calibration (SP-TAAC + SLR-WBC) is applied, it recovers near-identical high accuracies across all training scenarios ($\sim 43\%$ average accuracy).
\end{itemize}

This methodological study exposes the limits of training-time trajectory alignment in model merging and demonstrates that post-hoc representation alignment is highly robust, successfully recovering representations even from completely unconstrained, drifted experts.

\section{Related Work}

\subsection{Parameter-Space Model Merging}
Model merging aims to unify multiple specialized expert networks fine-tuned from a shared progenitor into a single multi-task model \cite{Choshen2022, Tang2023}. Simple Weight Averaging (WA) directly averages expert weights \cite{Wortsman2022}. Modern parameter-space techniques, such as TIES-Merging \cite{Yadav2023} and DARE \cite{Yu2023}, prune and align parameter updates to reduce task interference. Other approaches include permutation alignment to align representations prior to merging \cite{Ainsworth2023, Stoica2023} and weight scaling/regularization \cite{Jin2023a, Matena2022}. However, because these methods operate purely in parameter space, they do not adapt to activation-level dynamics at test time and still suffer from representation collapse.

\subsection{Activation Calibration and Alignment}
Rather than modifying weights, representation calibration operates in activation space. REPAIR \cite{Jordan2023} scales activations to match the expected variance of the original experts. Task-Conditional Activation Calibration (TCAC) \cite{Anonymous2024b} performs channel-wise affine scaling but requires an oracle task-ID at test time. SP-TAAC \cite{Anonymous2023} mitigates this via global, layer-wise scaling, while SLR-WBC \cite{Anonymous2024c} provides an analytical, closed-form weight and BatchNorm correction that achieves representation alignment with zero test-time inference overhead. However, existing works evaluate these calibration methods solely on top of simple weight averaging under unregularized expert fine-tuning, leaving a critical methodological gap.

\subsection{Regularized Fine-Tuning}
Regularization techniques are widely used in transfer learning to prevent catastrophic forgetting and overfitting. L2-SP (L2-Regularization targeting Starting Parameters) penalizes the L2 distance of active weights from their pre-trained starting weights, preserving pre-trained features \cite{Li2018}. In this work, we leverage L2-SP and weight decay as diagnostic tools to control expert parameter drift and study its direct impact on post-merge representation collapse.

\section{Methodological Framework}

We formalize our evaluation framework by defining our metrics for expert parameter drift, update alignment, and post-merge calibration.

\subsection{Parameter Drift and Update Alignment}
Let $W_{\text{init}} \in \mathbb{R}^D$ represent the flattened parameters of the shared pre-trained progenitor model. Let $W_k \in \mathbb{R}^D$ represent the parameters of the expert model fine-tuned on task $k \in \{1, \dots, K\}$. The parameter update (or task vector) of expert $k$ is defined as:
\begin{equation}
    \Delta W_k = W_k - W_{\text{init}}
\end{equation}
We measure the magnitude of expert parameter drift using the Euclidean L2 distance from initialization:
\begin{equation}
    \mathcal{D}(W_k) = \|W_k - W_{\text{init}}\|_2
\end{equation}
To analyze how independent fine-tuning trajectories coordinate in parameter space, we compute the cosine similarity between the update vectors of different tasks $i$ and $j$:
\begin{equation}
    \mathcal{S}(i, j) = \frac{\langle \Delta W_i, \Delta W_j \rangle}{\|\Delta W_i\|_2 \|\Delta W_j\|_2}
\end{equation}
A cosine similarity of $\sim 0$ indicates orthogonal trajectories, while higher values indicate coordinated update paths.

\subsection{Corrected SP-TAAC Calibration}
During our investigation, we identified a critical scale mismatch bug in standard SP-TAAC implementations. Prior formulations compared the standard deviation of merged activations before BatchNorm to expert activations after BatchNorm, causing severe scale distortions. In our corrected SP-TAAC calibration, we register hooks on the outputs of the Conv layers (before-BN) for both the merged and expert models.

For an early layer $l$ with merged Conv output activations $V \in \mathbb{R}^{B \times C \times H \times W}$ and expert Conv output activations $\{V_{\text{expert}, k}\}_{k=1}^K$, we compute the channel-wise standard deviations $\sigma_{\text{merged}}, \sigma_{\text{target}} \in \mathbb{R}^C$ over the joint calibration dataset $\mathcal{D}_{\text{cal}}$:
\begin{equation}
    \sigma_{\text{merged}, c} = \text{std}(V_{:, c, :, :})
\end{equation}
\begin{equation}
    \sigma_{\text{target}, c} = \text{std}([V_{\text{expert}, 1, :, c, :, :}, \dots, V_{\text{expert}, K, :, c, :, :}])
\end{equation}
The scale correction factor $\gamma \in \mathbb{R}^C$ is computed and clamped to $[0.1, 10.0]$:
\begin{equation}
    \gamma_c = \frac{\sigma_{\text{target}, c}}{\sigma_{\text{merged}, c} + \epsilon}
\end{equation}
We then scale the merged BatchNorm running weights and biases:
\begin{equation}
    w_l \leftarrow \gamma \cdot w_l, \quad b_l \leftarrow \gamma \cdot b_l
\end{equation}
By keeping both tensors strictly before-BN, our corrected SP-TAAC alignment completely avoids representational distortion.

\subsection{SLR-WBC Low-Rank Weight Correction}
For deeper layers where representation collapse is highly localized, we apply SLR-WBC weight and BatchNorm correction. For a convolutional layer, let $X \in \mathbb{R}^{d_{\text{in}} \times M}$ be the unfolded and flattened input activations. Let $W_{\text{curr}} \in \mathbb{R}^{C_{\text{out}} \times d_{\text{in}}}$ be the current Conv weights. The target Conv output $V_{\text{target}} \in \mathbb{R}^{C_{\text{out}} \times M}$ is obtained by inverting the experts' BatchNorm operations:
\begin{equation}
    V_{\text{target}, k, c} = \frac{H_{\text{target}, k, c} - b_{k, c}}{w_{k, c}} \sqrt{\sigma_{k, c}^2 + \epsilon} + \mu_{k, c}
\end{equation}
where $H_{\text{target}, k}$ is the expert BN output. We solve the uncentered ridge regression problem to find the optimal full-rank weight correction $\Delta W^*$:
\begin{equation}
    \Delta W^* = E X^T (X X^T + \lambda I)^{-1}
\end{equation}
where $E = V_{\text{target}} - W_{\text{curr}} X$ and $\lambda = \text{reg} \cdot M$. We obtain the rank-$r$ truncated Singular Value Decomposition (SVD) of $\Delta W^*$:
\begin{equation}
    \Delta W^{(r)} = U_{:, :r} \Sigma_{:r, :r} V^T_{:, :r}
\end{equation}
We update the Conv weights in-place $W_{\text{curr}} \leftarrow W_{\text{curr}} + \Delta W^{(r)}$, recompute the new activations $V_{\text{new}}$, and update the BatchNorm running statistics to their empirical mean and variance. Finally, we solve a channel-wise linear regression to align the BN outputs with $H_{\text{target}}$.

\subsection{Theoretical Analysis of Variance Collapse and SVD Regularization}
\label{sec:theoretical_analysis}
To ground our empirical investigations, we provide a formal mathematical treatment of two key phenomena: (1) why multi-task weight averaging triggers representation and variance collapse, and how trajectory alignment mitigates it; and (2) how ridge regularization stabilizes the SVD calibration in small-budget regimes.

\paragraph{Mathematical Analysis of Variance Collapse.}
Consider a single linear layer (or unfolded convolutional layer) in a merged backbone. Let $W_{\text{init}} \in \mathbb{R}^{C_{\text{out}} \times C_{\text{in}}}$ be the shared pre-trained progenitor, and let $W_k = W_{\text{init}} + \Delta W_k$ represent the weights of expert $k \in \{1, \dots, K\}$ fine-tuned on task $k$. The simple Weight Averaged (WA) weights are given by:
\begin{equation}
    W_{\text{merged}} = \frac{1}{K} \sum_{k=1}^K W_k = W_{\text{init}} + \frac{1}{K} \sum_{k=1}^K \Delta W_k
\end{equation}
Let $x \in \mathbb{R}^{C_{\text{in}}}$ be an input activation vector drawn from a distribution across tasks. The output representation of the merged layer is $y_{\text{merged}} = W_{\text{merged}} x$, whereas the output of expert $k$ is $y_k = W_k x = W_{\text{init}} x + \Delta W_k x$. 

To isolate the behavior of task-specific updates, let $z_{\text{init}} = W_{\text{init}} x$ and $z_k = \Delta W_k x$ be the pre-trained and task-specific activation components. We can express the output of expert $k$ as $y_k = z_{\text{init}} + z_k$. The multi-task target representation space is represented by the mixture of expert outputs, whose variance across channels and tasks is:
\begin{equation}
    \text{Var}(y_{\text{target}}) = \text{Var}(z_{\text{init}}) + \frac{1}{K} \sum_{k=1}^K \text{Var}(z_k)
\end{equation}
In contrast, the representation produced by the merged model is $y_{\text{merged}} = z_{\text{init}} + \frac{1}{K} \sum_{k=1}^K z_k$. The variance of this merged representation is:
\begin{equation}
    \text{Var}(y_{\text{merged}}) = \text{Var}(z_{\text{init}}) + \text{Var}\left(\frac{1}{K} \sum_{k=1}^K z_k\right)
\end{equation}
Assuming the task activation components $z_k$ have identical variance $\sigma^2$ and symmetric pairwise Pearson correlation coefficient $\rho \in [-1, 1]$, the variance of the averaged update component is:
\begin{equation}
    \text{Var}\left(\frac{1}{K} \sum_{k=1}^K z_k\right) = \frac{1 + (K - 1)\rho}{K} \sigma^2
\end{equation}
This formulation reveals the fundamental mechanics of representation collapse:
\begin{itemize}
    \item \textbf{Orthogonal Updates ($\rho \approx 0$)}: When experts are trained with zero or standard regularization, they drift in orthogonal directions. In this unconstrained regime, the task updates destructively interfere, reducing the update variance to $\frac{1}{K} \sigma^2$. For a large number of tasks $K$, this causes severe variance collapse in deep layers, destroying multi-task representational capacity.
    \item \textbf{Aligned Updates ($\rho > 0$)}: When training-time regularization is introduced (e.g., Scenario C), the expert update directions are constrained to a shared manifold, raising the correlation $\rho$. This directly increases the merged variance toward the target $\sigma^2$ (achieving full recovery at $\rho = 1$), explaining why high weight decay naturally suppresses representation collapse.
\end{itemize}

\paragraph{SVD Ridge Regularization and Generalization.}
In the post-merge calibration phase, we compute the low-rank weight correction $\Delta W^* = E X^T (X X^T + \lambda I)^{-1}$ over $M = N \times H \times W$ calibration activation columns. 

When the calibration budget $N$ is small, the activation matrix $X \in \mathbb{R}^{d_{\text{in}} \times M}$ contains very few samples, making the empirical covariance matrix $X X^T$ highly ill-conditioned or low-rank. In the unregularized limit ($\lambda \to 0$), the inverse $(X X^T)^{-1}$ exhibits extremely large eigenvalues, which amplifies any high-frequency noise or sample-specific artifacts in the error matrix $E$. This causes the weight correction $\Delta W^*$ to overfit to the small calibration sample, leading to poor generalization on the test set.

By adding the scale-dependent ridge regularization $\lambda = \text{reg} \cdot M$, we structurally bound the eigenvalues of the regularized inverse:
\begin{equation}
    \lambda_{\max}\left( (X X^T + \lambda I)^{-1} \right) \le \frac{1}{\lambda} = \frac{1}{\text{reg} \cdot M}
\end{equation}
This ridge term prevents arbitrary amplification of the weight corrections along low-variance activation directions, ensuring that the computed low-rank update $\Delta W^{(r)}$ is highly stable and robustly generalizable to the test set, even when estimated from as few as $N=16$ samples.

\section{Experimental Evaluation}

\subsection{Experimental Setup}
We evaluate our framework on MNIST, Fashion-MNIST, and CIFAR-10. Following \cite{Anonymous2024c}, we fine-tune ImageNet-pretrained ResNet-18 expert models on $5,000$-sample subsets of each dataset for $5$ epochs. Grayscale images are resized to $32\times32$ and replicated to $3$ channels. The experts are trained using SGD (momentum $0.9$, learning rate $1\times10^{-4}$, dropout $0.1$). Calibration uses $N=128$ samples per task, and evaluation is conducted on the full $10,000$-sample test sets. We evaluate four distinct training-time regularization scenarios:
\begin{itemize}\itemsep=0pt \parsep=0pt
    \item \textbf{Scenario A (Low Reg)}: Weight decay = $0.0$.
    \item \textbf{Scenario B (Std Reg)}: Weight decay = $1\times10^{-4}$.
    \item \textbf{Scenario C (High Reg)}: Weight decay = $1\times10^{-2}$.
    \item \textbf{Scenario D (L2-SP Reg)}: Explicit L2 penalty on distance from progenitor with $\lambda = 1\times10^{-3}$.
\end{itemize}

\subsection{Expert Drift and Update Alignment Results}
We analyze how training-time regularization affects the spatial configuration of the expert weights. The results are summarized in \cref{tab:drift_metrics} and visualized in \cref{fig:weight_alignment}.

\begin{table}[t]
\caption{Expert Parameter Drift and Weight Update Alignment across Regularization Scenarios}
\label{tab:drift_metrics}
\vskip 0.15in
\begin{center}
\begin{small}
\resizebox{\columnwidth}{!}{%
\begin{tabular}{lcccc}
\toprule
Metric & Scenario A (Low Reg) & Scenario B (Std Reg) & Scenario C (High Reg) & Scenario D (L2-SP Reg) \\
\midrule
MNIST Weight L2-Drift & 0.7668 & 0.7762 & 0.8285 & 0.7667 \\
F-MNIST Weight L2-Drift & 0.7149 & 0.7141 & 0.7866 & 0.7141 \\
CIFAR-10 Weight L2-Drift & 0.6769 & 0.6772 & 0.7492 & 0.6769 \\
\midrule
Cos Sim: MNIST / F-MNIST & 0.0431 & 0.0414 & \textbf{0.2037} & 0.0434 \\
Cos Sim: MNIST / CIFAR-10 & -0.0023 & -0.0018 & \textbf{0.1869} & 0.0046 \\
Cos Sim: F-MNIST / CIFAR-10 & -0.0088 & -0.0109 & \textbf{0.1900} & -0.0123 \\
\bottomrule
\end{tabular}%
}
\end{small}
\end{center}
\vskip -0.1in
\end{table}

\begin{figure*}[t]
    \centering
    \begin{subfigure}[b]{0.48\textwidth}
        \centering
        \includegraphics[width=\textwidth,height=0.18\textheight]{weight_alignment.png}
        \caption{Expert Weight Update Alignment}
        \label{fig:weight_alignment}
    \end{subfigure}
    \hfill
    \begin{subfigure}[b]{0.48\textwidth}
        \centering
        \includegraphics[width=\textwidth,height=0.18\textheight]{merged_accuracies.png}
        \caption{Multi-Task Merged Model Average Test Accuracy}
        \label{fig:merged_accuracies}
    \end{subfigure}
    \caption{Expert weight alignment and merged accuracy. (a) High Weight Decay (Scenario C) coordinates fine-tuning trajectories, raising update cosine similarity to over $0.19$. (b) Our Hybrid calibration consistently improves average accuracy by $+17.4\%$ absolute, acting as a robust equalizer.}
    \label{fig:alignment_and_accuracies}
    \vspace{-10pt}
\end{figure*}

Our analysis reveals several critical insights:
\begin{enumerate}
    \item \textbf{Trajectory Coordination}: Under standard training (Scenario A and B) or mild L2-SP (Scenario D), the expert update directions are completely orthogonal, with cosine similarities hovering around $0.00$ to $0.04$. However, applying High Weight Decay (Scenario C) coordinates the independent training paths, raising cross-task cosine similarity to over $\textbf{0.19}$.
    \item \textbf{Constrained Drift}: While weight decay is traditionally viewed as a simple L2 norm constraint on weights, in pre-trained models it acts as a soft anchor that regularizes updates back toward the progenitor initialization manifold, enforcing shared directionality across different classification tasks.
\end{enumerate}

\subsection{Multi-Task Merging and Calibration Results}
We merge the experts from each scenario using Weight Averaging and evaluate them under three calibration settings. The results are summarized in \cref{tab:merged_acc} and visualized in \cref{fig:merged_accuracies}.

\begin{table}[t]
\caption{Multi-Task Merged Model Average Test Accuracies (\%) across Calibration Settings}
\label{tab:merged_acc}
\vskip 0.15in
\begin{center}
\begin{small}
\resizebox{\columnwidth}{!}{%
\begin{tabular}{lcccc}
\toprule
Calibration Method & Scenario A (Low Reg) & Scenario B (Std Reg) & Scenario C (High Reg) & Scenario D (L2-SP Reg) \\
\midrule
Uncalibrated (none) & 26.20\% & 27.26\% & \textbf{27.88\%} & 27.30\% \\
SP-TAAC Calibration & 35.36\% & 35.95\% & 35.54\% & \textbf{36.22\%} \\
Hybrid (Ours, $r=4$) & \textbf{43.63\%} & 43.10\% & 43.27\% & 42.96\% \\
\bottomrule
\end{tabular}%
}
\end{small}
\end{center}
\vskip -0.1in
\end{table}

Our comparative analysis highlights several crucial methodological findings:
\begin{enumerate}\itemsep=0pt \parsep=0pt
    \item \textbf{Regularization Benefits the Base Merge}: For the uncalibrated merged model (`none`), increasing training-time regularization steadily improves accuracy (from **$26.20\%$** in Scenario A to **$27.88\%$** in Scenario C). This confirms our hypothesis that constraining expert parameter drift directly suppresses the severity of representation collapse.
    \item \textbf{The Post-Hoc Equalizer Phenomenon}: Despite the clear benefits of trajectory alignment on the uncalibrated base model, the post-merge calibration methods act as a powerful equalizer. When our Hybrid calibration (SP-TAAC + SLR-WBC) is applied, it achieves near-identical high accuracies across all scenarios (**$43.63\%$** in Scenario A vs **$43.27\%$** in Scenario C).
    \item \textbf{Efficacy of Representation Alignment}: Our Hybrid calibration provides a massive absolute accuracy improvement of over **$+17.43\%$** compared to the uncalibrated baseline. This highlights that while training-time alignment is beneficial, post-hoc representation alignment is highly robust, successfully recovering representational capacity even from highly drifted, unregularized experts.
\end{enumerate}

\subsection{Hyperparameter Sensitivity and SVD Diagnostics}
To establish an even deeper methodological understanding of representation alignment, we perform a thorough hyperparameter grid search over SVD rank $r \in \{1, 2, 4, 8\}$ and ridge regularization strength $\text{reg} \in \{0.1, 0.5, 1.0\}$. We evaluate these configurations across Scenario A (Low Reg) and Scenario C (High Reg). 

As shown in \cref{tab:grid_search_results}, increasing SVD rank $r$ dramatically improves performance. At rank $r=8$ and a scale-invariant regularization of $\text{reg}=0.1$, the Hybrid calibration reaches a stellar **$48.43\%$** average accuracy (Scenario A) and **$48.35\%$** (Scenario C), representing a massive **$+22.23\%$** absolute increase over the uncalibrated model. Crucially, the final performance under proper calibration remains virtually identical regardless of whether the models were trained with standard, unregularized parameters (Scenario A) or high regularization (Scenario C), further solidifying the ``post-hoc equalizer'' thesis.

\begin{table}[h]
\vspace{-5pt}
\caption{SLR-WBC Hyperparameter Grid Search (Average Accuracy \%) across Ranks and Regularizations}
\label{tab:grid_search_results}
\vskip 0.1in
\begin{center}
\begin{small}
\resizebox{\columnwidth}{!}{%
\begin{tabular}{lcccccccc}
\toprule
& \multicolumn{2}{c}{Rank 1} & \multicolumn{2}{c}{Rank 2} & \multicolumn{2}{c}{Rank 4} & \multicolumn{2}{c}{Rank 8} \\
\cmidrule(lr){2-3} \cmidrule(lr){4-5} \cmidrule(lr){6-7} \cmidrule(lr){8-9}
Reg ($\text{reg}$) & Scen A & Scen C & Scen A & Scen C & Scen A & Scen C & Scen A & Scen C \\
\midrule
reg 0.1 & 37.50\% & 36.15\% & 43.05\% & 42.26\% & 45.61\% & 45.33\% & \textbf{48.43\%} & \textbf{48.35\%} \\
reg 0.5 & 38.49\% & 38.07\% & 41.62\% & 40.90\% & 43.63\% & 43.27\% & 44.78\% & 44.84\% \\
reg 1.0 & 37.58\% & 38.01\% & 40.89\% & 39.88\% & 43.00\% & 42.28\% & 43.38\% & 43.25\% \\
\bottomrule
\end{tabular}%
}
\end{small}
\end{center}
\vskip -0.1in
\vspace{-10pt}
\end{table}

Furthermore, we extract and profile activations from the deepest block's final BatchNorm layer (\texttt{layer4.1.bn2}) to directly diagnose representation collapse. 

\begin{figure*}[t]
    \centering
    \begin{subfigure}[b]{0.48\textwidth}
        \centering
        \includegraphics[width=\textwidth,height=0.18\textheight]{variance_collapse_recovery.png}
        \caption{Activation Standard Deviation across Channels}
        \label{fig:variance_collapse_recovery}
    \end{subfigure}
    \hfill
    \begin{subfigure}[b]{0.48\textwidth}
        \centering
        \includegraphics[width=\textwidth,height=0.18\textheight]{spectrum_decay_recovery.png}
        \caption{Feature Spectrum Decay (SVD) at \texttt{layer4.1.bn2}}
        \label{fig:spectrum_decay_recovery}
    \end{subfigure}
    \caption{Diagnostics at \texttt{layer4.1.bn2}. (a) The uncalibrated model exhibits severe variance collapse, which our Hybrid calibration fully restores to the Oracle target. (b) The uncalibrated model experiences rapid singular value decay (loss of rank), whereas our Hybrid calibration successfully recovers high-dimensional capacity.}
    \label{fig:diagnostics_recovery}
    \vspace{-10pt}
\end{figure*}

In \cref{fig:variance_collapse_recovery}, we analyze the standard deviation across channels. The uncalibrated model exhibits severe variance collapse, rendering activations extremely low in magnitude and heavily compressed. Our Hybrid calibration completely aligns and restores these standard deviations, matching the Oracle experts' profile.

In \cref{fig:spectrum_decay_recovery}, we evaluate the normalized singular values (SVD feature spectrum). The uncalibrated merged model experiences a rapid, exponential decay of singular values, indicating severe dimensional collapse (the representational capacity is restricted to a very small subspace). Our Hybrid calibration successfully recovers the high-dimensional representational capacity, aligning perfectly with the Oracle spectrum.

\subsection{The Calibration Data Budget and Regularization Sweep}
Activation calibration relies on a small budget $N$ of calibration samples. Extremely small budgets ($N \le 32$) risk representational overfitting, while larger budgets increase VRAM and inference costs. To analyze this trade-off, we sweep budget $N \in \{16, 32, 64, 128, 256\}$ and SVD regularization parameter $\text{reg} \in \{0.01, 0.1, 0.5, 1.0, 5.0\}$ using the Hybrid method (at rank $r=4$) across Scenario A and Scenario C experts.

The quantitative results are visualized in \cref{fig:calibration_budget_analysis} and analyzed below.

\begin{figure*}[t]
    \centering
    \includegraphics[width=0.96\textwidth,height=0.19\textheight]{calibration_budget_analysis.png}
    \caption{Hybrid calibration performance across calibration sizes $N$ and regularization strengths $\text{reg}$. Stronger ridge regularization ($\text{reg} \ge 1.0$) stabilizes small budgets ($N=16$), preventing overfitting and restoring test performance.}
    \label{fig:calibration_budget_analysis}
    \vspace{-10pt}
\end{figure*}

Our sweep yields several profound methodological insights:
\begin{enumerate}\itemsep=0pt \parsep=0pt \topsep=0pt
    \item \textbf{Overfitting in Small-Budget Regimes}: When the calibration budget is extremely restricted ($N=16$), standard low-regularization settings ($\text{reg}=0.01$) exhibit a noticeable drop in test performance (averaging around $38.5\%$). This performance drop confirms our hypothesis of representation overfitting, where the low-rank correction $\Delta W^{(r)}$ over-aligns the merged model to the specific activation coordinates of the $16$ calibration samples, degrading generalizability.
    \item \textbf{Ridge Regularization as a Stabilizer}: Strikingly, increasing the SVD ridge regularization parameter $\text{reg}$ completely mitigates this overfitting. In the $N=16$ regime, raising $\text{reg}$ from $0.01$ to $1.0$ or $5.0$ restores the merged model's average accuracy to over $42\%$, effectively bridging the generalization gap. Mathematically, a larger $\text{reg}$ dampens the magnitude of the full-rank weight update $\Delta W^* = E X^T (X X^T + \lambda I)^{-1}$ by heavily penalizing large correction steps, thereby ensuring that the rank-$r$ projection remains highly conservative and generalizable.
    \item \textbf{Asymptotic Convergence}: As the calibration budget expands to $N \ge 64$, the performance across all regularization levels converges asymptotically to the optimal $\sim 43\%$ range. This demonstrates that with sufficient calibration data, the empirical covariance matrix $X X^T$ is naturally well-conditioned and robust, reducing the model's sensitivity to the specific value of the ridge regularization parameter.
    \item \textbf{Cross-Scenario Invariance}: Crucially, the qualitative and quantitative behavior of the budget-regularization interaction is identical across both Scenario A (Low Regularization) and Scenario C (High Regularization) experts. This cross-scenario invariance further validates the "post-hoc equalizer" thesis: the post-merge representation alignment is so robust that it completely governs the final model performance, irrespective of whether the experts were regularized during training.
\end{enumerate}

\subsection{Multi-Seed Robustness Evaluation}
To ensure statistical significance, we run our training, merging, and calibration pipeline across three independent random seeds ($\text{seed} \in \{42, 43, 44\}$). We report means and standard deviations for all core metrics to confirm our findings are robust to random variations.

Our multi-seed analysis yields the following robust conclusions:
\begin{itemize}\itemsep=0pt \parsep=0pt
    \item \textbf{Trajectory Alignment Robustness}: For Scenario A (Low Reg), the average expert update cosine similarity across the three seeds is $0.018 \pm 0.007$, representing almost perfectly orthogonal trajectories. In contrast, for Scenario C (High Reg), the average similarity is $0.200 \pm 0.006$. This confirms that training-time regularization consistently and robustly coordinates expert trajectories in parameter space across random initializations.
    \item \textbf{Uncalibrated Merged Model Accuracies}: Under Scenario A, the uncalibrated merged model achieves an average test accuracy of $26.28\% \pm 1.13\%$ across seeds. Under Scenario C, it achieves $26.64\% \pm 1.40\%$. Regularization consistently helps the base merge, albeit by a small margin.
    \item \textbf{The Equalizer Invariance}: Strikingly, under our optimal Hybrid calibration ($r=8, \text{reg}=0.1$), the average merged accuracy reaches $47.00\% \pm 1.26\%$ for Scenario A and $47.33\% \pm 0.93\%$ for Scenario C. This identical performance (within standard error) across seeds completely solidifies our "post-hoc equalizer" thesis: activation-space calibration acts as a robust equalizer that completely dominates final performance, regardless of whether experts were regularized during training or what random seed was used.
\end{itemize}

\subsection{Task Vector Scaling Baseline Audit}
In Task Vector Scaling (or Task Arithmetic), task-specific updates are scaled by $\lambda \in (0, 1]$ before merging:
\begin{equation}
    W_{\text{merged}} = W_{\text{init}} + \lambda \sum_{k=1}^K (W_k - W_{\text{init}})
\end{equation}
We sweep $\lambda \in [0.1, 1.0]$ under Scenario A and Scenario C to test if simple scaling can bypass post-hoc activation calibration.

Our baseline audit reveals several critical methodological insights:
\begin{enumerate}\itemsep=0pt \parsep=0pt \topsep=0pt
    \item \textbf{Uncalibrated Scaling Brittleness}: Without calibration, the merged model is highly sensitive to the scale of $\lambda$. It peaks at $\lambda = 0.3$ (near the simple Weight Averaging scale of $1/K = 0.33$), achieving $27.70\%$ (Scenario A) and $28.77\%$ (Scenario C) average accuracy. However, scaling slightly further to $\lambda = 0.4$ drops the performance instantly to a random-guessing limit of $9.93\%$, showing that the uncalibrated model experiences a catastrophic, abrupt representation collapse cliff.
    \item \textbf{Calibration-Enabled Peak Performance}: Strikingly, when our Hybrid calibration ($r=8, \text{reg}=0.1$) is applied, it not only stabilizes the model at $\lambda = 0.4$ but recovers an absolute peak multi-task performance of $\textbf{49.88\%}$ (Scenario A) and $\textbf{50.41\%}$ (Scenario C)! This is $+2.0\%$ higher than standard Weight Averaging.
    \item \textbf{Catastrophic Instability Cliff}: For $\lambda \ge 0.5$, both the uncalibrated and calibrated models collapse. Under these extreme scales, the activation covariance matrices become mathematically singular (with zero diagonal elements), causing the SVD ridge solver to fail due to complete representation collapse. This reveals the physical limits of post-hoc calibration: while it is exceptionally powerful at correcting moderate scaling distortions, it cannot recover representations when parameters are scaled past the stability manifold.
\end{enumerate}

\subsection{TIES-Merging and DARE-Merging Baseline Audit}
We extend our methodological audit to modern parameter-space pruning techniques: TIES-Merging \cite{Yadav2023} and DARE-Merging \cite{Yu2023}. We merge the Scenario A (Low Reg) and Scenario C (High Reg) experts using TIES (pruning rate $p=20\%$) and DARE (drop rate $p=80\%$), evaluated uncalibrated and with our optimal Hybrid calibration ($r=8, \text{reg}=0.1$).

The quantitative results are summarized in \cref{tab:ties_dare_comparison}. They demonstrate that pure parameter-space techniques actually exacerbate uncalibrated representation collapse. Without calibration, TIES-Merging drops to $\sim 14.5\%$ accuracy, while DARE completely collapses to random guessing ($9.93\%$). However, under our Hybrid post-hoc calibration, TIES-Merging is fully resurrected, reaching a peak accuracy of \textbf{49.37\%} (Scenario C). In contrast, DARE remains collapsed because its random pruning destroys representation spaces so severely that activation covariance matrices become mathematically singular (precluding least-squares calibration). This reinforces the "post-hoc equalizer" thesis: post-hoc alignment is highly robust and synergizes beautifully with parameter-space trimming, but extreme random dropout destroys the physical manifold beyond recovery.

\begin{table}[h]
\vspace{-5pt}
\caption{TIES and DARE Average Test Accuracies (\%)}
\label{tab:ties_dare_comparison}
\vskip 0.1in
\begin{center}
\begin{small}
\resizebox{\columnwidth}{!}{%
\begin{tabular}{lcccc}
\toprule
& \multicolumn{2}{c}{Scenario A (Low Reg)} & \multicolumn{2}{c}{Scenario C (High Reg)} \\
\cmidrule(lr){2-3} \cmidrule(lr){4-5}
Merge Method & Uncalibrated & Hybrid Calibrated & Uncalibrated & Hybrid Calibrated \\
\midrule
Weight Averaging (WA) & 26.20\% & 48.43\% & 27.88\% & 48.35\% \\
TIES-Merging & 14.55\% & \textbf{48.82\%} & 14.05\% & \textbf{49.37\%} \\
DARE-Merging & 9.93\% & Failed (Singular) & 9.93\% & Failed (Singular) \\
\bottomrule
\end{tabular}%
}
\end{small}
\end{center}
\vskip -0.1in
\vspace{-10pt}
\end{table}

\section{Methodological Discussion}
From \textit{The Methodologist}'s perspective, these findings reveal a critical, unexamined assumption in the model merging literature. Prior works argue that representation collapse is an inevitable physical consequence of parameter interpolation, necessitating increasingly complex test-time activation calibration. Our work deconstructs this assumption by showing that:
\begin{itemize}\itemsep=0pt \parsep=0pt
    \item Representation collapse is a direct symptom of unregularized training trajectories. By introducing a modest high weight decay during expert training, practitioners can double update alignment and naturally suppress representation collapse in parameter space.
    \item However, the benefits of training-time alignment do not stack with post-hoc calibration. Because post-hoc calibration is mathematically designed to align activation scales and manifolds directly, it acts as a robust equalizer, rendering the initial parameter-space drift largely irrelevant once calibration is complete.
\end{itemize}
These insights suggest that future research should focus less on extremely complex training-time regularization to aid merging, and instead prioritize robust, zero-overhead post-hoc calibration methods like SLR-WBC that are highly resilient to expert training-time configurations.

\section{Conclusion}
In this paper, we presented a rigorous methodological deconstruction of representation collapse in multi-task model merging. By evaluating independent ResNet-18 experts trained across four regularization levels and merged under three calibration configurations, we exposed the critical relationship between training-time parameter drift and post-merge representation alignment. We demonstrated that while training-time regularization coordinates training trajectories (raising update cosine similarity from $0.04$ to $0.20$) and improves the uncalibrated merged model, post-hoc calibration remains an essential and robust equalizer, consistently delivering an outstanding $+22.23\%$ absolute accuracy boost. Our findings establish a new, fairer benchmarking standard for the model merging community.

\bibliography{example_paper}
\bibliographystyle{icml2026}

\newpage
\appendix
\onecolumn
\icmltitle{Appendix: The Class Bias Confounder in Post-Hoc Calibration}

\section{Calibration Generalization under Class-Skewed Calibration Sets}
\label{appendix:class_bias}

\subsection{Motivation and Setup}
Activation calibration techniques, such as Sparsity-Preserving Task-Agnostic Activation Calibration (SP-TAAC) and SVD-based Low-Rank Weight and BatchNorm Calibration (SLR-WBC), are fundamentally grounded in aligning activation statistics over a small calibration dataset $\mathcal{D}_{\text{cal}}$. Existing literature implicitly assumes that $\mathcal{D}_{\text{cal}}$ is a clean, fully representative sample drawn from the exact same multi-task distribution as the test set. However, in practical deployment scenarios, this assumption is frequently violated: calibration datasets may suffer from severe class imbalances or complete class omissions due to localized data collection constraints.

To investigate this realistic challenge under our critical \textit{Methodologist} lens, we design an extremely rigorous generalization audit. We evaluate our Hybrid post-hoc calibration ($r=8$) under two distinct calibration dataset configurations:
\begin{itemize}
    \item \textbf{Unbiased Calibration}: Standard, uniformly selected samples ($N=128$) across all 10 classes of MNIST, Fashion-MNIST, and CIFAR-10.
    \item \textbf{Skewed Calibration}: A biased calibration set ($N=128$) where 50\% of the classes are completely omitted during the calibration phase. Specifically, for MNIST and Fashion-MNIST, the calibration set only contains even-numbered classes ($\{0, 2, 4, 6, 8\}$). For CIFAR-10, it only contains the first 5 classes ($\{0, 1, 2, 3, 4\}$).
\end{itemize}

During evaluation, we partition the full $10,000$-sample test set into two groups:
\begin{enumerate}
    \item \textbf{In-Distribution (ID) Classes}: The classes present in the skewed calibration set (e.g., even classes).
    \item \textbf{Out-of-Distribution (OOD) Classes}: The classes completely omitted from the skewed calibration set (e.g., odd classes).
\end{enumerate}
We evaluate the performance of both Scenario A (Low Reg) and Scenario C (High Reg) merged backbones. We sweep the SVD ridge regularization parameter $\text{reg} \in \{0.1, 1.0, 5.0\}$ under skewed calibration to see if stronger regularization can stabilize representation space alignment.

\subsection{Quantitative Results and Analysis}
The complete quantitative breakdown is presented in \cref{tab:class_bias_results}. Our empirical findings expose a major, previously unexamined vulnerability in post-merge activation calibration.

\begin{table*}[h]
\caption{Generalization performance (ID / OOD / Total Accuracy \%) of post-hoc Hybrid calibration under Unbiased vs. Class-Skewed calibration sets.}
\label{tab:class_bias_results}
\vskip 0.15in
\begin{center}
\begin{small}
\resizebox{\linewidth}{!}{%
\begin{tabular}{llcccccc}
\toprule
& & \multicolumn{3}{c}{\textbf{Scenario A (Low Reg)}} & \multicolumn{3}{c}{\textbf{Scenario C (High Reg)}} \\
\cmidrule(lr){3-5} \cmidrule(lr){6-8}
Task & Calibration Config & ID Acc (\%) & OOD Acc (\%) & Total Acc (\%) & ID Acc (\%) & OOD Acc (\%) & Total Acc (\%) \\
\midrule
\textbf{MNIST} & Uncalibrated (none) & 29.01\% & 21.19\% & 25.04\% & 33.84\% & 22.29\% & 27.98\% \\
& Unbiased ($\text{reg}=0.1$) & 60.43\% & 69.69\% & 65.13\% & 58.14\% & 70.08\% & 64.20\% \\
& Unbiased ($\text{reg}=1.0$) & 59.64\% & 63.50\% & 61.60\% & 55.87\% & 63.05\% & 59.51\% \\
& Skewed ($\text{reg}=0.1$) & \textbf{77.10\%} & 11.79\% & 43.96\% & \textbf{73.45\%} & 15.08\% & 43.83\% \\
& Skewed ($\text{reg}=1.0$) & 69.61\% & 10.78\% & 39.76\% & 65.92\% & 13.91\% & 39.53\% \\
& Skewed ($\text{reg}=5.0$) & 67.19\% & 5.62\% & 35.95\% & 62.12\% & 8.89\% & 35.11\% \\
\midrule
\textbf{F-MNIST} & Uncalibrated (none) & 37.92\% & 27.34\% & 32.63\% & 40.48\% & 29.32\% & 34.90\% \\
& Unbiased ($\text{reg}=0.1$) & 44.26\% & 68.00\% & 56.13\% & 42.92\% & 70.36\% & 56.64\% \\
& Unbiased ($\text{reg}=1.0$) & 39.72\% & 48.00\% & 43.86\% & 39.24\% & 51.48\% & 45.36\% \\
& Skewed ($\text{reg}=0.1$) & \textbf{50.06\%} & 0.96\% & 25.51\% & \textbf{49.02\%} & 1.68\% & 25.35\% \\
& Skewed ($\text{reg}=1.0$) & 43.04\% & 0.48\% & 21.76\% & 40.76\% & 0.58\% & 20.67\% \\
& Skewed ($\text{reg}=5.0$) & 39.62\% & 0.32\% & 19.97\% & 38.24\% & 0.28\% & 19.26\% \\
\midrule
\textbf{CIFAR-10} & Uncalibrated (none) & 24.70\% & 17.16\% & 20.93\% & 23.70\% & 17.80\% & 20.75\% \\
& Unbiased ($\text{reg}=0.1$) & 27.58\% & 20.50\% & 24.04\% & 28.40\% & 20.00\% & 24.20\% \\
& Unbiased ($\text{reg}=1.0$) & 29.02\% & 20.36\% & 24.69\% & 29.30\% & 20.44\% & 24.87\% \\
& Skewed ($\text{reg}=0.1$) & 25.56\% & 18.96\% & 22.26\% & 27.78\% & 18.60\% & 23.19\% \\
& Skewed ($\text{reg}=1.0$) & 28.50\% & 18.44\% & 23.47\% & \textbf{30.34\%} & 17.74\% & 24.04\% \\
& Skewed ($\text{reg}=5.0$) & 28.34\% & 18.84\% & 23.59\% & 30.42\% & 17.60\% & 24.01\% \\
\bottomrule
\end{tabular}%
}
\end{small}
\end{center}
\vskip -0.1in
\end{table*}

\subsection{Deep Methodological Insights}
\begin{enumerate}
    \item \textbf{The Overfitting Catastrophe}: Under skewed calibration, the model's accuracy on the seen (ID) classes rises significantly compared to the unbiased calibration setup. For example, on MNIST (Scenario A) under skewed calibration with $\text{reg}=0.1$, ID accuracy reaches a high of \textbf{77.10\%} (vs. 60.43\% unbiased). However, the accuracy on unseen (OOD) classes catastrophically collapses to a dismal \textbf{11.79\%} (vs. 69.69\% unbiased). On Fashion-MNIST, the collapse is even more severe, with OOD accuracy dropping to an almost-zero \textbf{0.96\%} (vs. 68.00\% unbiased). This empirical result exposes a fundamental, unexamined vulnerability of activation calibration: it aligns activation statistics purely to fit the manifolds of the classes seen in the calibration set, in the process distorting and corrupting the representations of unseen classes.
    \item \textbf{The Inadequacy of Regularization as a Defense}: Increasing the SVD ridge regularization parameter $\text{reg}$ from $0.1$ to $5.0$ reduces the overfitting on ID classes (e.g., MNIST ID accuracy drops from $77.10\%$ to $67.19\%$), but it does not rescue the collapsed OOD performance. On Fashion-MNIST, OOD accuracy remains below $0.5\%$. This indicates that the representation distortion caused by class-skewed activation statistics cannot be simply regularized away; the calibration process fundamentally corrupts the shared features of unseen classes by forcing them to conform to the statistics of a biased class subset.
    \item \textbf{Cross-Scenario Invariance}: Crucially, the behavior is identical across both Scenario A (Low Regularization) and Scenario C (High Regularization) experts. This cross-scenario invariance solidifies our "post-hoc equalizer" thesis under a new dimension: training-time trajectory regularization is entirely incapable of protecting the model against the distribution-shift vulnerabilities of post-hoc calibration.
\end{enumerate}

These findings indicate that future research must prioritize developing class-robust post-hoc calibration methods or explicitly auditing calibration sets for representation bias before deployment, establishing a new and essential direction for evaluation protocols in the model merging field.

\end{document}
